\textbf{COLUNI 2021 -- Questão 26}\par\medskip A astrônoma Cecilia Payne-Gaposchkin (1900 – 1979), aos 25 anos, defendeu sua tese de doutorado intitulada \textit{Stellar Atmospheres}. Depois de analisar milhares de espectros estelares gravados em placas de vidro e utilizando a Teoria da Ionização, concluiu que as atmosferas das estrelas consistiam, predominantemente, de hidrogênio e hélio. Na ocasião, a teoria vigente ditava que a composição da atmosfera das estrelas era semelhante à da Terra. Cecilia foi, então, convencida, pelos seus supervisores (Harlow Shapley e Henry Norris Russell), a considerar os seus dados possíveis de erro. Quatro anos depois, por outros métodos, Russell chegou aos mesmos resultados e os publicou, tendo recebido os créditos da descoberta. Posteriormente, a tese de Cecilia Payne foi considerada a mais brilhante tese de astronomia! Mesmo com doutorado completo, trabalhou muitos anos recebendo um salário ínfimo, por ser mulher, e, somente em 1956, foi nomeada professora em Harvard. Fonte: FABBIANO, G. The woman who explained the stars. \textbf{Nature}, vol. 578, p. 509-510, 2020. (Adaptado)

\par

A composição química da atmosfera do Sol é, em percentual de massa, 71\% de hidrogênio, 27\% de hélio e o restante de outros elementos químicos, entre eles oxigênio, carbono e ferro. Disponível em: http://hyperphysics.phy-astr.gsu.edu/hbase/Tables/suncomp.html. Acesso em: 20 nov. 2020. (Adaptado)

\par

Com base nas informações fornecidas e nos conhecimentos sobre ciências, assinale a afirmativa CORRETA:\begin{enumerate}
\item O número de prótons de um átomo de deutério, um dos isótopos do hidrogênio, difere em uma unidade do
        número de prótons do átomo de hélio.
\item Considerando a massa do Sol igual a $2{,}0 \times 10^{30}$ kg e a composição química da atmosfera do Sol,
        apresentada no enunciado da questão, a massa de hélio, nessa estrela, corresponde a $5{,}4 \times 10^{30}$ g.
\item A massa atômica do elemento químico hélio é o dobro da massa atômica do elemento químico hidrogênio.
\item O elemento químico hélio está localizado na décima oitava coluna da Tabela Periódica, sendo classificado
        como um halogênio.\end{enumerate}
